% Source: source/普通高考/2012/2012天津理.pdf, page 1, question 3.
% pdfimages -list reports no embedded bitmap; the program flowchart is vector/text artwork.
% Grid evidence: tmp/redraw_flowchart_2012_tianjin_science/grid/q3_grid.jpg.
% Comparison crop: tmp/redraw_flowchart_2012_tianjin_science/grid/q3_original_clean.png,
% taken from the ungridded page render with a safety border.
% Nodes: 开始; 输入 x; |x|>1?; x=√|x|−1; x=2x+1; 输出 x; 结束.
% Directed edges: 开始→输入→判定; 判定(是)→平方根更新→左侧回环→判定;
% 判定(否)→x=2x+1→输出→结束. The loop and output branches are independent.
\begin{tikzpicture}[
  >=Stealth,
  line width=0.55pt,
  line cap=round,
  line join=round,
  every node/.style={font=\normalsize,anchor=center,align=center,
    execute at begin node={\setlength{\parindent}{0pt}\noindent}},
  flowbox/.style={draw,minimum height=.68cm,inner xsep=4pt,inner ysep=2pt,
    anchor=center,align=center},
  terminal/.style={flowbox,rounded corners=.18cm,minimum width=1.35cm,
    minimum height=.72cm},
  io/.style={flowbox,shape=trapezium,trapezium left angle=72,
    trapezium right angle=108,minimum height=.78cm},
  decision/.style={flowbox,shape=diamond,aspect=2.85,inner sep=2pt},
  flowarrow/.style={->,line width=0.55pt},
]
  \node[terminal] (start) at (0,9.15) {开始};
  \node[io,minimum width=2.35cm] (input) at (0,7.82) {输入 $x$};
  \coordinate (loopjoin) at (0,6.82);
  \node[decision,minimum width=3.55cm,minimum height=1.02cm] (test) at (0,5.98)
    {$|x|>1?$};
  \node[flowbox,minimum width=3.65cm,minimum height=.98cm] (sqrtupdate) at (0,4.42)
    {$x=\sqrt{|x|}-1$};
  \node[flowbox,minimum width=2.70cm] (affine) at (0,2.72)
    {$x=2x+1$};
  \node[io,minimum width=1.95cm] (output) at (0,1.36) {输出 $x$};
  \node[terminal] (finish) at (0,0.06) {结束};

  \draw[flowarrow] (start.south) -- (input.north);
  \draw (input.south) -- (loopjoin);
  \draw[flowarrow] (loopjoin) -- (test.north);
  \draw[flowarrow] (test.south) -- node[right=2pt] {是} (sqrtupdate.north);
  % Affirmative branch loops left into the stem above the decision.
  \coordinate (loop-left) at (-2.45,3.92);
  \coordinate (loop-left-top) at (-2.45,6.82);
  \draw (sqrtupdate.south) -- (loop-left) -- (loop-left-top);
  \draw[flowarrow] (loop-left-top) -- (loopjoin);

  % Negative branch bypasses the loop and enters the affine update.
  \coordinate (out-right) at (3.20,5.98);
  \coordinate (out-right-bottom) at (3.20,3.42);
  \coordinate (affine-entry) at (0,3.42);
  \draw (test.east) -- (out-right) -- (out-right-bottom);
  \draw (out-right-bottom) -- (affine-entry);
  \draw[flowarrow] (affine-entry) -- (affine.north);
  \draw[flowarrow] (affine.south) -- (output.north);
  \draw[flowarrow] (output.south) -- (finish.north);
  \node[anchor=south west,inner sep=0pt] at (1.95,6.10) {否};

  \path[use as bounding box] (-2.85,-0.25) rectangle (3.85,9.60);
\end{tikzpicture}
